\begin{enumerate}[label=\thesubsection.\arabic*, ref=\thesubsection.\theenumi, itemsep=0pt, leftmargin=0pt, labelwidth=0pt, itemindent=0pt, labelsep=1em, align=right]

\item Rank of a Matrix
\begin{align}
\rank\brak{\vec{A}} &= \text{number of non-zero rows of } \vec{A} \text{ in row echelon form}
\label{eq:F-rank-pivots}
\end{align}

\item $p$-Series Test
\begin{align}
\sum_{n=1}^{\infty} \frac{1}{n^{p}} \text{ converges} &\iff p > 1
\label{eq:F-pseries}
\end{align}

\item Telescoping Series
\begin{align}
\sum_{n=1}^{N} \frac{1}{n\brak{n+1}} &= \sum_{n=1}^{N} \brak{\frac{1}{n} - \frac{1}{n+1}}
\\
&= 1 - \frac{1}{N+1}
\label{eq:F-telescope}
\end{align}

\item Exponential Series
\begin{align}
e^x &= \sum_{n=0}^{\infty} \frac{x^n}{n!}
\\
\sum_{n=1}^{\infty} \frac{1}{n!} &= e - 1
\label{eq:F-exp}
\end{align}
\solution Substituting $x = 1$ and removing the term $n = 0$ gives \eqref{eq:F-exp}.

\item Ratio Test For a series $\sum a_n$,
\begin{align}
L &= \lim_{n \to \infty} \abs{\frac{a_{n+1}}{a_n}}
\\
L < 1 &\implies \text{the series converges}
\label{eq:F-ratio}
\end{align}

\item Expectation
\begin{align}
E\sbrak{X} &= \sum_{x} x \, \pr{X = x}
\\
E\sbrak{X_1 + X_2} &= E\sbrak{X_1} + E\sbrak{X_2}
\label{eq:F-expect}
\end{align}

\item Marginal PMF
\begin{align}
\pr{X = x} &= \sum_{y} \pr{X = x, Y = y}
\label{eq:F-marginal}
\end{align}

\item Independent Random Variables If $X$ and $Y$ are independent,
\begin{align}
\pr{X = x, Y = y} &= \pr{X = x}\pr{Y = y}
\label{eq:F-indep}
\end{align}

\item PMF of a Sum of Independent Random Variables
\begin{align}
p_{X+Y}\sbrak{k} &= \sum_{i} p_X\sbrak{i}\, p_Y\sbrak{k - i}
\label{eq:F-conv-pmf}
\end{align}
\solution The event $\cbrak{X + Y = k}$ is the union over $i$ of the disjoint events $\cbrak{X = i, Y = k - i}$. Hence
\begin{align}
\pr{X + Y = k} &= \sum_{i} \pr{X = i, Y = k - i}
\\
&= \sum_{i} \pr{X = i}\pr{Y = k - i}
\\
&= \sum_{i} p_X\sbrak{i}\, p_Y\sbrak{k - i}
\end{align}
upon using \eqref{eq:F-indep}.

\item Uncorrelated Random Variables For independent, zero-mean, unit-variance $a_n$ and $a_m$, $n \ne m$,
\begin{align}
E\sbrak{a_n a_m} &= E\sbrak{a_n}E\sbrak{a_m} = 0
\\
E\sbrak{a_n^2} &= 1
\\
E\sbrak{a_n a_m} &= \delta_{n - m}
\label{eq:F-uncorr}
\end{align}

\item Gaussian Symmetry The density of a zero-mean Gaussian random variable $N$ is symmetric, so
\begin{align}
\pr{N > c} &= \pr{N < -c}
\label{eq:F-gauss-sym}
\end{align}

\item Admittances
\begin{align}
Y_R &= \frac{1}{R}
\\
Y_L &= \frac{1}{\jj \omega L}
\\
Y_C &= \jj \omega C
\label{eq:F-admittance}
\end{align}

\item Ohm's Law for Phasors
\begin{align}
\underline{I} &= Y \underline{V}
\label{eq:F-phasor-ohm}
\end{align}

\item Polar Form of Phasors
\begin{align}
\jj &= 1 \angle 90^\circ
\\
-\jj &= 1 \angle -90^\circ
\\
\brak{a \angle \theta_1}\brak{b \angle \theta_2} &= ab \angle \brak{\theta_1 + \theta_2}
\label{eq:F-polar}
\end{align}

\item Series $RLC$ Branch
\begin{align}
\abs{Z} &= \sqrt{R^2 + \brak{\omega L - \frac{1}{\omega C}}^2}
\label{eq:F-series-rlc}
\\
\omega L &= \frac{1}{\omega C}
\\
C &= \frac{1}{\omega^2 L}
\label{eq:F-resonance}
\end{align}

\item $Z$-Parameters of a Two-Port
\begin{align}
\myvec{V_1 \\ V_2} &= \myvec{z_{11} & z_{12} \\ z_{21} & z_{22}} \myvec{I_1 \\ I_2}
\label{eq:F-zdef}
\\
z_{11} &= \frac{V_1}{I_1}\Big|_{I_2 = 0}
\\
z_{21} &= \frac{V_2}{I_1}\Big|_{I_2 = 0}
\label{eq:F-zparam}
\end{align}

\item Parallel Resistors
\begin{align}
R_1 \parallel R_2 &= \frac{R_1 R_2}{R_1 + R_2}
\label{eq:F-rnet}
\end{align}

\item LTI Systems For a system $y\sbrak{n} = T\cbrak{x\sbrak{n}}$,
\begin{align}
\delta\sbrak{n} * h\sbrak{n} &= h\sbrak{n}
\label{eq:F-conv}
\\
T\cbrak{a x_1 + b x_2} &= a T\cbrak{x_1} + b T\cbrak{x_2} \quad \text{(linear)}
\label{eq:F-linear}
\\
T\cbrak{0} &= 0 \quad \text{(a linear system)}
\label{eq:F-zero}
\\
T\cbrak{x\sbrak{n - k}} &= y\sbrak{n - k} \quad \text{(time invariant)}
\label{eq:F-ti}
\end{align}

\item Fourier Transform
\begin{align}
F\brak{\omega} &= \int_{-\infty}^{\infty} f\brak{t} e^{-\jj \omega t} \dd t
\label{eq:F-ft}
\end{align}

\item Magnitude Bound of the Fourier Transform
\begin{align}
\abs{F\brak{\omega}} &\le \int_{-\infty}^{\infty} \abs{f\brak{t}} \dd t
\label{eq:F-ft-bound}
\end{align}
\solution Using $\abs{\int g} \le \int \abs{g}$ and $\abs{e^{-\jj \omega t}} = 1$ in \eqref{eq:F-ft},
\begin{align}
\abs{F\brak{\omega}} &\le \int_{-\infty}^{\infty} \abs{f\brak{t}} \abs{e^{-\jj \omega t}} \dd t
\\
&= \int_{-\infty}^{\infty} \abs{f\brak{t}} \dd t
\end{align}

\item Transform of $e^{-\abs{t}}$
\begin{align}
e^{-\abs{t}} &\longleftrightarrow \frac{2}{1 + \omega^2}
\label{eq:F-ft-pair}
\end{align}
\solution From \eqref{eq:F-ft},
\begin{align}
\int_{-\infty}^{\infty} e^{-\abs{t}} e^{-\jj \omega t} \dd t &= \int_{-\infty}^{0} e^{t\brak{1 - \jj \omega}} \dd t + \int_{0}^{\infty} e^{-t\brak{1 + \jj \omega}} \dd t
\\
&= \frac{1}{1 - \jj \omega} + \frac{1}{1 + \jj \omega}
\\
&= \frac{2}{1 + \omega^2}
\end{align}

\item Fourier Series and Time Scaling A periodic signal with fundamental period $T_0$ has
\begin{align}
f\brak{t} &= \sum_{k = -\infty}^{\infty} c_k e^{\jj 2\pi k t / T_0}
\label{eq:F-fs}
\end{align}
and, for $\alpha > 0$,
\begin{align}
f\brak{\alpha t} &= \sum_{k = -\infty}^{\infty} c_k e^{\jj 2\pi k \alpha t / T_0}
\label{eq:F-fs-scale}
\\
T_y &= \frac{T_0}{\alpha}
\label{eq:F-fs-period}
\end{align}
\solution Replacing $t$ by $\alpha t$ in \eqref{eq:F-fs} yields \eqref{eq:F-fs-scale}, with the same coefficients $c_k$. For $y\brak{t} = f\brak{\alpha t}$,
\begin{align}
y\brak{t + \frac{T_0}{\alpha}} &= f\brak{\alpha t + T_0}
\\
&= f\brak{\alpha t} = y\brak{t}
\end{align}
yielding \eqref{eq:F-fs-period}.

\item Impulse Sampling With $\omega_s = \frac{2\pi}{T_s}$,
\begin{align}
\sum_{n} \delta\brak{t - \tau - n T_s} &= \frac{1}{T_s} \sum_{k} e^{\jj k \omega_s \brak{t - \tau}}
\label{eq:F-comb-fs}
\\
Y\brak{\omega} &= \frac{1}{T_s} \sum_{k} e^{-\jj k \omega_s \tau} F\brak{\omega - k \omega_s}
\label{eq:F-sampled}
\end{align}
where $y\brak{t} = f\brak{t} \sum_{n} \delta\brak{t - \tau - n T_s}$.
\solution Multiplying $f\brak{t}$ by \eqref{eq:F-comb-fs} and using that $f\brak{t}e^{\jj k \omega_s t}$ has transform $F\brak{\omega - k\omega_s}$,
\begin{align}
y\brak{t} &= \frac{1}{T_s} \sum_{k} e^{-\jj k \omega_s \tau} f\brak{t} e^{\jj k \omega_s t}
\\
Y\brak{\omega} &= \frac{1}{T_s} \sum_{k} e^{-\jj k \omega_s \tau} F\brak{\omega - k \omega_s}
\end{align}

\item Ideal Lowpass Reconstruction
\begin{align}
H\brak{\omega} &= T_s, \quad \abs{\omega} < \omega_c
\label{eq:F-lpf}
\\
\omega_s > 2\omega_c &\iff T_s < \frac{\pi}{\omega_c}
\label{eq:F-nyquist}
\end{align}

\item $z$-Transform and DTFT
\begin{align}
X\brak{z} &= \sum_{n} x\sbrak{n} z^{-n}
\\
X\brak{e^{\jj \omega}} &= \sum_{n} x\sbrak{n} e^{-\jj \omega n}
\label{eq:F-ztrans}
\end{align}
The DTFT exists if the region of convergence (ROC) contains the unit circle. The ROC contains no poles but may contain zeros. For $x\sbrak{n} = \alpha \delta\sbrak{n}$, $X\brak{z} = \alpha$ with the ROC the entire $z$-plane.

\item Characteristic Equation of a Unity-Feedback Loop
\begin{align}
1 + L\brak{s} &= 0
\label{eq:F-char}
\end{align}

\item Stability of a Quadratic Polynomial
\begin{align}
a_2 s^2 + a_1 s + a_0 \text{ is stable} &\iff a_2, a_1, a_0 \text{ have the same sign}
\label{eq:F-hurwitz2}
\end{align}

\item Roots of $s^3 + c$ For $c > 0$ and $k = 0, 1, 2$,
\begin{align}
s^3 + c = 0 &\implies s = c^{1/3} e^{\jj \brak{\pi + 2\pi k}/3}
\label{eq:F-cubic-roots}
\end{align}

\item Root Locus A point $s_0$ lies on the root locus of $1 + K G\brak{s} = 0$ for some $K \ge 0$ iff
\begin{align}
G\brak{s_0} &= -\frac{1}{K}
\label{eq:F-rootlocus}
\end{align}
i.e.\ $G\brak{s_0}$ is real and negative.
\solution At a closed-loop pole $1 + K G\brak{s_0} = 0$, hence
\begin{align}
G\brak{s_0} &= -\frac{1}{K}
\\
K &= \frac{1}{\abs{G\brak{s_0}}}
\end{align}

\item Mason's Gain Formula
\begin{align}
T &= \frac{\sum_{k} P_k \Delta_k}{\Delta}
\label{eq:F-mason}
\\
\Delta &= 1 - \sum_{i} L_i + \cdots
\end{align}
where $P_k$ are the forward-path gains, $L_i$ the loop gains and $\Delta_k$ is $\Delta$ with the loops touching path $k$ removed.

\item State Equation For $\dot{x}_i = \lambda_i x_i + b_i u$,
\begin{align}
x_i\brak{t} &= e^{\lambda_i t} x_i\brak{0} + b_i \int_{0}^{t} e^{\lambda_i \brak{t - \sigma}} u\brak{\sigma} \dd \sigma
\label{eq:F-state}
\\
\abs{x_i\brak{t}} &\le \abs{x_i\brak{0}} + \frac{\abs{b_i} U}{\abs{\lambda_i}} \quad \text{if } \lambda_i < 0, \abs{u} \le U
\label{eq:F-state-bounded}
\\
x_i\brak{t} &= \frac{b_i}{\lambda_i}\brak{e^{\lambda_i t} - 1} \quad \text{if } \lambda_i > 0, u \equiv 1, x_i\brak{0} = 0
\label{eq:F-state-unbounded}
\\
x_i\brak{t} &= b_i t \quad \text{if } \lambda_i = 0, u \equiv 1, x_i\brak{0} = 0
\label{eq:F-state-integrator}
\end{align}
\solution Putting $u \equiv 1$ and $x_i\brak{0} = 0$ in \eqref{eq:F-state},
\begin{align}
x_i\brak{t} &= b_i \int_{0}^{t} e^{\lambda_i \brak{t - \sigma}} \dd \sigma
\\
&= \frac{b_i}{\lambda_i}\brak{e^{\lambda_i t} - 1}, \quad \lambda_i \ne 0
\\
&= b_i t, \quad \lambda_i = 0
\end{align}
which grows without bound for $\lambda_i \ge 0$ and $b_i \ne 0$.

\item Majority Function
\begin{align}
Y_{\mathrm{maj}} &= AB + AC + BC
\label{eq:F-maj}
\end{align}

\item XOR with a Constant
\begin{align}
Y \oplus 0 &= Y
\\
Y \oplus 1 &= \overline{Y}
\label{eq:F-xor}
\end{align}

\item Two's Complement Subtraction
\begin{align}
X + \overline{Y} + 1 &\equiv X - Y \pmod{2^n}
\label{eq:F-twos}
\end{align}
\solution For $n$-bit numbers $\overline{Y} = \brak{2^n - 1} - Y$, so
\begin{align}
X + \overline{Y} + 1 &= X + \brak{2^n - 1} - Y + 1
\\
&= X - Y + 2^n
\\
&\equiv X - Y \pmod{2^n}
\end{align}

\item (D Flip-Flop and Setup Constraint)
\begin{align}
Q^{+} &= D
\label{eq:F-dff}
\\
T_{\mathrm{clk}} &\ge t_{cq} + t_{comb} + t_{setup}
\\
f_{\max} &= \frac{1}{T_{\mathrm{clk,min}}}
\label{eq:F-setup}
\end{align}

\item Frequency Modulation
\begin{align}
f_i\brak{t} &= \frac{1}{2\pi} \frac{d\phi\brak{t}}{dt}
\\
\Delta f_{\mathrm{peak}} &= \max_{t} \abs{f_i\brak{t} - f_c}
\label{eq:F-fm}
\end{align}

\item Envelope Detection
\begin{align}
\text{envelope (DSB-SC)} &= A_c \abs{m\brak{t}}
\\
\text{envelope (AM)} &= A_c \abs{1 + \mu m\brak{t}}
\label{eq:F-envelope}
\end{align}

\item Maximum-Likelihood Detection For equal priors and equal-variance Gaussian noise, the detector chooses the nearest symbol
\begin{align}
\hat{S} &= \arg\min_{s} \abs{y - s}
\label{eq:F-ml}
\end{align}

\item Ideal ADC
\begin{align}
\mathrm{SNR}_{\max} &= 6.02 N + 1.76 \text{ dB}
\\
R_b &= f_s N
\label{eq:F-adc}
\end{align}

\item Inverting Op-Amp Filter
\begin{align}
H\brak{s} &= -\frac{Z_f\brak{s}}{Z_i\brak{s}}
\label{eq:F-inv}
\end{align}
A factor $\frac{1}{1 + sRC}$ gives a corner frequency
\begin{align}
\omega &= \frac{1}{RC}
\label{eq:F-corner}
\end{align}

\item Entropy
\begin{align}
H\brak{X} &= -\sum_{x} p\brak{x} \log_2 p\brak{x}
\label{eq:F-entropy}
\end{align}

\item Mutual Information
\begin{align}
I\brak{X;Y} &= \sum_{x,y} p\brak{x,y} \log_2 \frac{p\brak{x,y}}{p\brak{x}p\brak{y}}
\label{eq:F-mi}
\\
I\brak{X;Y} &= 0 \quad \text{if } X, Y \text{ are independent}
\label{eq:F-mi-zero}
\end{align}
\solution For independent $X, Y$, $p\brak{x,y} = p\brak{x}p\brak{y}$ by \eqref{eq:F-indep}, so every logarithm in \eqref{eq:F-mi} is $\log_2 1 = 0$:
\begin{align}
I\brak{X;Y} &= \sum_{x,y} p\brak{x,y} \log_2 1
\\
&= 0
\end{align}

\item Minimum Distance of a Binary Linear Code With $w$ the Hamming weight,
\begin{align}
d_{\min} &= \min_{\vec{c} \ne \vec{0}} w\brak{\vec{c}}
\label{eq:F-dmin}
\end{align}

\item Faraday's Law
\begin{align}
\abs{\mathcal{E}} &= \abs{\frac{d\Phi_B}{dt}} = A \abs{\frac{dB}{dt}}
\\
I &= \frac{\abs{\mathcal{E}}}{R_{\mathrm{total}}}
\label{eq:F-faraday}
\end{align}

\item Current Density
\begin{align}
J &= \sigma E
\\
I &= J A
\label{eq:F-J}
\end{align}

\item Surface Charge
\begin{align}
Q &= \iint \rho_s \dd A
\label{eq:F-charge}
\end{align}

\item Ideal Diode and Phase-to-Time Conversion An ideal diode is forward biased iff the anode voltage exceeds the cathode voltage. A phase interval $\Delta \theta$ at angular frequency $\omega$ lasts
\begin{align}
\Delta t &= \frac{\Delta \theta}{\omega}
\label{eq:F-angletime}
\end{align}

\item Second Derivative Test
\begin{align}
f'\brak{x_0} = 0, f''\brak{x_0} > 0 &\implies \text{local minimum}
\label{eq:F-second}
\\
f'\brak{x_0} = 0, f''\brak{x_0} < 0 &\implies \text{local maximum}
\end{align}

\item Bounds on an Integral
\begin{align}
m \le f\brak{x} \le M \text{ on } \sbrak{a, b} &\implies m\brak{b - a} \le \int_{a}^{b} f\brak{x} \dd x \le M\brak{b - a}
\label{eq:F-intbound}
\end{align}

\item Euler--Cauchy Equation For $t^2 y'' + a t y' + b y = 0$ with $y = t^m$,
\begin{align}
m\brak{m - 1} + a m + b &= 0
\label{eq:F-euler-cauchy}
\end{align}
\solution Substituting $y = t^m$, $y' = m t^{m-1}$ and $y'' = m\brak{m-1}t^{m-2}$,
\begin{align}
t^2 y'' + a t y' + b y &= t^m \sbrak{m\brak{m - 1} + a m + b}
\\
&= 0
\end{align}
yielding \eqref{eq:F-euler-cauchy}.

\item Least Squares The scalar $x$ minimizing $\norm{x\vec{a} - \vec{b}}$ is
\begin{align}
x^* &= \frac{\vec{a}^\top \vec{b}}{\vec{a}^\top \vec{a}}
\label{eq:F-ls}
\end{align}
\solution Let $J\brak{x} = \norm{x\vec{a} - \vec{b}}^2$. Then
\begin{align}
J\brak{x} &= x^2 \vec{a}^\top \vec{a} - 2x \vec{a}^\top \vec{b} + \vec{b}^\top \vec{b}
\\
\frac{dJ}{dx} &= 2x \vec{a}^\top \vec{a} - 2\vec{a}^\top \vec{b} = 0
\\
x^* &= \frac{\vec{a}^\top \vec{b}}{\vec{a}^\top \vec{a}}
\end{align}

\item Cauchy's Theorem For a closed counter-clockwise contour $C$,
\begin{align}
\oint_{C} f\brak{z} \dd z &= 0 \quad \text{if } f \text{ is analytic inside and on } C
\label{eq:F-cauchy}
\end{align}

\item Integral of $z^n$ on the Unit Circle
\begin{align}
\oint_{\abs{z} = 1} z^n \dd z &= \begin{cases} 2\pi \jj, & n = -1 \\ 0, & n \ne -1 \end{cases}
\label{eq:F-zn}
\end{align}
\solution Putting $z = e^{\jj \theta}$ and $dz = \jj e^{\jj \theta} d\theta$,
\begin{align}
\oint_{\abs{z} = 1} z^n \dd z &= \jj \int_{0}^{2\pi} e^{\jj \brak{n + 1}\theta} \dd \theta
\\
&= \begin{cases} 2\pi \jj, & n = -1 \\ 0, & n \ne -1 \end{cases}
\end{align}

\item Poles of $\sec z$ The poles of $\sec z = \frac{1}{\cos z}$ are where $\cos z = 0$:
\begin{align}
z &= \frac{\pi}{2} + k\pi, \quad k \in \mathbb{Z}
\label{eq:F-sec}
\end{align}

\end{enumerate}
